
\subsection{Correction}

\begin{correctionbox}[Exercice 1 :]
\begin{enumerate}
    \item Pour que $f(x)$ soit une PDF valide, on doit avoir $\int_{-\infty}^{\infty} f(x)\,dx = 1$. Alors :
    \[
    \int_0^2 c x^2 \, dx = c \left[ \frac{x^3}{3} \right]_0^2 = c \cdot \frac{8}{3} = 1 \implies c = \frac{3}{8}.
    \]
    
    \item Avec $c = \frac{3}{8}$, on calcule :
    \[
    P(X > 1) = \int_1^2 \frac{3}{8} x^2 \, dx = \frac{3}{8} \left[ \frac{x^3}{3} \right]_1^2 = \frac{3}{8} \left( \frac{8}{3} - \frac{1}{3} \right) = \frac{3}{8} \cdot \frac{7}{3} = \frac{7}{8}.
    \]
\end{enumerate}
\end{correctionbox}

\begin{correctionbox}[Exercice 2 :]
\begin{enumerate}
    \item La CDF est définie par $F(x) = \int_{-\infty}^x f(t)\,dt$.
    \item Si $x < 0$, $F(x) = 0$.
        \item Si $x \in [0,2]$, $F(x) = \int_0^x \frac{1}{2} t \, dt = \frac{1}{2} \cdot \frac{t^2}{2} \Big|_0^x = \frac{x^2}{4}$.
        \item Si $x > 2$, $F(x) = 1$.
    Donc :
    \[
    F(x) = 
    \begin{cases}
    0 & x < 0 \\
    \frac{x^2}{4} & 0 \le x \le 2 \\
    1 & x > 2
    \end{cases}
    \]

    \item $P(1 \le X \le 1.5) = F(1.5) - F(1) = \frac{(1.5)^2}{4} - \frac{1^2}{4} = \frac{2.25 - 1}{4} = \frac{1.25}{4} = 0.3125$.
\end{enumerate}
\end{correctionbox}

\begin{correctionbox}[Exercice 3 :]
La PDF est la dérivée de la CDF :
\[
f(x) = F'(x) = 
\begin{cases}
0 & x < 0 \\
\cos(x) & 0 \le x \le \pi/2 \\
0 & x > \pi/2
\end{cases}
\]
Donc $f(x) = \cos(x)$ pour $x \in [0, \pi/2]$, et $0$ sinon.

\end{correctionbox}

\begin{correctionbox}[Exercice 4 :]
Soit $X$ continue avec PDF $f(x)$. Alors :
\[
P(X = a) = \int_a^a f(x)\,dx = 0,
\]
car l'intégrale sur un intervalle de longueur nulle est nulle.

\end{correctionbox}

\begin{correctionbox}[Exercice 5 :]
\begin{enumerate}
    \item $E[X] = \int_0^1 x \cdot 3x^2 \, dx = 3 \int_0^1 x^3 \, dx = 3 \cdot \frac{1}{4} = \frac{3}{4}$.
    
    \item $E[X^2] = \int_0^1 x^2 \cdot 3x^2 \, dx = 3 \int_0^1 x^4 \, dx = 3 \cdot \frac{1}{5} = \frac{3}{5}$.
    
    \item $\text{Var}(X) = E[X^2] - (E[X])^2 = \frac{3}{5} - \left(\frac{3}{4}\right)^2 = \frac{3}{5} - \frac{9}{16} = \frac{48 - 45}{80} = \frac{3}{80}$.
\end{enumerate}
\end{correctionbox}

\begin{correctionbox}[Exercice 6 :]
Par LOTUS : $E[1/X] = \int_0^1 \frac{1}{x} \cdot 2x \, dx = \int_0^1 2 \, dx = 2$.
Mais attention : en $x=0$, $1/x$ n'est pas défini. Toutefois, comme il s'agit d'un point de mesure nulle, l'intégrale converge et vaut 2.

\end{correctionbox}

\begin{correctionbox}[Exercice 7 :]
On utilise $\text{Var}(Y) = E[Y^2] - (E[Y])^2$ avec $Y = aX + b$ :
\[
E[aX + b] = aE[X] + b, \quad E[(aX + b)^2] = E[a^2X^2 + 2abX + b^2] = a^2E[X^2] + 2abE[X] + b^2.
\]
Alors :
\begin{align*}
\text{Var}(aX + b) &= E[(aX + b)^2] - (E[aX + b])^2 \\
&= a^2E[X^2] + 2abE[X] + b^2 - (aE[X] + b)^2 \\
&= a^2E[X^2] + 2abE[X] + b^2 - \left(a^2(E[X])^2 + 2abE[X] + b^2\right) \\
&= a^2E[X^2] - a^2(E[X])^2 = a^2 \text{Var}(X).
\end{align*}
\end{correctionbox}

\begin{correctionbox}[Exercice 8 :]
$T \sim \text{Unif}(0, 12)$, donc $f_T(t) = \frac{1}{12}$ pour $t \in [0,12]$.
\begin{enumerate}
    \item $P(T < 4) = \int_0^4 \frac{1}{12} dt = \frac{4}{12} = \frac{1}{3}$.
    \item $P(5 \le T \le 10) = \int_5^{10} \frac{1}{12} dt = \frac{5}{12}$.
\end{enumerate}
\end{correctionbox}

\begin{correctionbox}[Exercice 9 :]
\begin{enumerate}
    \item $E[T] = \frac{0 + 12}{2} = 6$ minutes.
    \item $\text{Var}(T) = \frac{(12 - 0)^2}{12} = \frac{144}{12} = 12$.
\end{enumerate}
\end{correctionbox}

\begin{correctionbox}[Exercice 10 :]
\[
E[X] = \int_a^b x \cdot \frac{1}{b-a} dx = \frac{1}{b-a} \int_a^b x \, dx = \frac{1}{b-a} \cdot \frac{b^2 - a^2}{2} = \frac{b^2 - a^2}{2(b-a)} = \frac{(b-a)(b+a)}{2(b-a)} = \frac{a+b}{2}.
\]
\end{correctionbox}

\begin{correctionbox}[Exercice 11 :]
On sait : $E[X] = \frac{a+b}{2} = 10$, $\text{Var}(X) = \frac{(b-a)^2}{12} = 12$.
De la première équation : $a + b = 20$.
De la seconde : $(b - a)^2 = 144 \implies b - a = 12$ (puisque $b > a$).
En résolvant le système :
\[
a + b = 20, \quad b - a = 12 \implies 2b = 32 \implies b = 16, \quad a = 4.
\]
Donc $a = 4$, $b = 16$.

\end{correctionbox}

\begin{correctionbox}[Exercice 12 :]
\begin{enumerate}
    \item $E[T] = \frac{1}{\lambda} = 5 \implies \lambda = \frac{1}{5} = 0.2$.
    \item $P(T \le 2) = 1 - e^{-\lambda \cdot 2} = 1 - e^{-0.4} \approx 1 - 0.6703 = 0.3297$.
\end{enumerate}
\end{correctionbox}

\begin{correctionbox}[Exercice 13 :]
Par la propriété d'absence de mémoire :
\[
P(T > 13 \mid T > 3) = P(T > 10) = e^{-\lambda \cdot 10} = e^{-2} \approx 0.1353.
\]
Et $P(T > 10) = e^{-2} \approx 0.1353$. Les deux probabilités sont égales.

\end{correctionbox}

\begin{correctionbox}[Exercice 14 :]
\[
E[X] = \int_0^\infty x \lambda e^{-\lambda x} dx.
\]
Intégration par parties : $u = x$, $dv = \lambda e^{-\lambda x} dx$, donc $du = dx$, $v = -e^{-\lambda x}$.
\[
E[X] = \left[ -x e^{-\lambda x} \right]_0^\infty + \int_0^\infty e^{-\lambda x} dx = 0 + \left[ -\frac{1}{\lambda} e^{-\lambda x} \right]_0^\infty = \frac{1}{\lambda}.
\]
\end{correctionbox}

\begin{correctionbox}[Exercice 15 :]
\[
P(X > s+t \mid X > s) = \frac{P(X > s+t)}{P(X > s)} = \frac{e^{-\lambda(s+t)}}{e^{-\lambda s}} = e^{-\lambda t} = P(X > t).
\]
\end{correctionbox}

\begin{correctionbox}[Exercice 16 :]
Soit $A = \{X \le a\}$, $B = \{X \le b\}$ avec $a < b$. Alors $A \subset B$, et :
\[
P(a < X \le b) = P(B \setminus A) = P(B) - P(A) = F(b) - F(a).
\]
\end{correctionbox}

\begin{correctionbox}[Exercice 17 :]
\[
E[aX + b] = \int_{-\infty}^\infty (ax + b) f_X(x) dx = a \int_{-\infty}^\infty x f_X(x) dx + b \int_{-\infty}^\infty f_X(x) dx = aE[X] + b.
\]
\end{correctionbox}

\begin{correctionbox}[Exercice 18 :]
$X \sim \text{Unif}(0,1)$, $Y = 5X - 2$. Transformation linéaire croissante.
CDF : $F_Y(y) = P(Y \le y) = P(5X - 2 \le y) = P\left(X \le \frac{y+2}{5}\right) = F_X\left(\frac{y+2}{5}\right)$.
Comme $X \in [0,1]$, alors $Y \in [-2, 3]$.
PDF : $f_Y(y) = f_X\left(\frac{y+2}{5}\right) \cdot \left|\frac{d}{dy} \left(\frac{y+2}{5}\right)\right| = 1 \cdot \frac{1}{5} = \frac{1}{5}$ pour $y \in [-2, 3]$.
Donc $Y \sim \text{Unif}(-2, 3)$.

\end{correctionbox}

\begin{correctionbox}[Exercice 19 :]
$X \sim \text{Unif}(-1,1)$, $f_X(x) = \frac{1}{2}$ sur $[-1,1]$, $Y = X^2$.
Pour $y \in [0,1]$, on a deux branches : $x = \sqrt{y}$ et $x = -\sqrt{y}$.
Formule générale pour transformation non monotone :
\[
f_Y(y) = \sum_{i} f_X(x_i) \left| \frac{dx_i}{dy} \right|, \quad \frac{dx}{dy} = \frac{d}{dy} (\pm \sqrt{y}) = \pm \frac{1}{2\sqrt{y}}.
\]
Donc :
\[
f_Y(y) = f_X(\sqrt{y}) \cdot \frac{1}{2\sqrt{y}} + f_X(-\sqrt{y}) \cdot \frac{1}{2\sqrt{y}} = \frac{1}{2} \cdot \frac{1}{2\sqrt{y}} + \frac{1}{2} \cdot \frac{1}{2\sqrt{y}} = \frac{1}{4\sqrt{y}} + \frac{1}{4\sqrt{y}} = \frac{1}{2\sqrt{y}}.
\]
Ainsi, $f_Y(y) = \frac{1}{2\sqrt{y}}$ pour $y \in (0,1]$, et $0$ sinon.

\end{correctionbox}

\begin{correctionbox}[Exercice 20 :]
$X \sim \text{Exp}(1)$, $f_X(x) = e^{-x}$ pour $x \ge 0$, $Y = \sqrt{X} \Rightarrow X = Y^2$.
Transformation strictement croissante pour $y \ge 0$.
\[
f_Y(y) = f_X(y^2) \cdot \left| \frac{d}{dy}(y^2) \right| = e^{-y^2} \cdot 2y, \quad y \ge 0.
\]
Donc $f_Y(y) = 2y e^{-y^2}$ pour $y \ge 0$.

\end{correctionbox}
